\documentclass[tikz,border=10pt]{standalone}
\usepackage{fontspec}
\usepackage{xeCJK}
\setmainfont{Times New Roman}
\setCJKmainfont{Hiragino Sans GB}
\setCJKsansfont{Microsoft YaHei}
\setCJKmonofont{STSong}
\usepackage{tikz}
\usetikzlibrary{matrix}

\tikzset{
  cell/.style={draw=black!70, line width=1.0pt, minimum width=4.2cm, minimum height=1.15cm, align=center, text width=3.8cm, fill=none},
  head/.style={cell, draw=blue!70!black, line width=1.2pt},
  left/.style={cell, draw=green!50!black, line width=1.2pt}
}

\begin{document}
\begin{tikzpicture}
  \matrix[matrix of nodes, nodes in empty cells, row sep=-\pgflinewidth, column sep=-\pgflinewidth] {
    \node[head] {项目}; &
    \node[head] {例 2}; &
    \node[head] {例 3}; \\
    \node[left] {开环形式}; &
    \node[cell] {$\frac{K^\ast}{s(s+1)(s+2)}$}; &
    \node[cell] {$K^\ast\frac{s+2}{s(s+1)(s+4)}$}; \\
    \node[left] {实轴根轨迹}; &
    \node[cell] {$(-\infty,-2]\cup[-1,0]$}; &
    \node[cell] {$[-4,-2]\cup[-1,0]$}; \\
    \node[left] {渐近线}; &
    \node[cell] {$\sigma_a=-1$, $\psi_a=\pm60^\circ,180^\circ$}; &
    \node[cell] {$\sigma_a=-\frac{3}{2}$, $\psi_a=\pm90^\circ$}; \\
    \node[left] {下一步}; &
    \node[cell] {求分离点}; &
    \node[cell] {求分离点与虚轴交点}; \\
  };
\end{tikzpicture}
\end{document}
